% !TeX root = ../thesis.tex


%The first chapter of your thesis.
\chapter{Changes w.r.t. the default template}\label{ch:introduction}



% Placeholder prose for previewing the campus layout; replace with your own text.
\section{Example text}

Lorem ipsum dolor sit amet, consectetur adipiscing elit. Sed do eiusmod tempor
incididunt ut labore et dolore magna aliqua. Ut enim ad minim veniam, quis
nostrud exercitation ullamco laboris nisi ut aliquip ex ea commodo consequat.
Duis aute irure dolor in reprehenderit in voluptate velit esse cillum dolore eu
fugiat nulla pariatur. Excepteur sint occaecat cupidatat non proident, sunt in
culpa qui officia deserunt mollit anim id est laborum.

Praesent vitae arcu non ipsum consequat tincidunt. Integer euismod, neque sed
finibus tincidunt, lorem erat malesuada nibh, vitae facilisis justo nulla in
urna. Aliquam erat volutpat. Donec tristique lacus vel sapien suscipit, sit
amet consequat neque interdum. Curabitur blandit tempus porttitor. Maecenas
faucibus mollis interdum. Vestibulum ante ipsum primis in faucibus orci luctus
et ultrices posuere cubilia curae; aenean eu leo quam. Pellentesque ornare sem
lacinia quam venenatis vestibulum.

Nullam quis risus eget urna mollis ornare vel eu leo. Cras mattis consectetur
purus sit amet fermentum. Donec ullamcorper nulla non metus auctor fringilla.
Vivamus sagittis lacus vel augue laoreet rutrum faucibus dolor auctor. Morbi
leo risus, porta ac consectetur ac, vestibulum at eros. Sed posuere consectetur
est at lobortis. Aenean lacinia bibendum nulla sed consectetur. Etiam porta sem
malesuada magna mollis euismod. Fusce dapibus, tellus ac cursus commodo, tortor
mauris condimentum nibh, ut fermentum massa justo sit amet risus.

Suspendisse potenti. Phasellus dignissim, sapien quis volutpat tincidunt, eros
nisl scelerisque turpis, nec tincidunt augue justo vitae neque. Nam eget erat
at magna pellentesque vulputate. Mauris non lorem sed augue malesuada
ullamcorper. Sed viverra nibh ut felis consequat, at scelerisque dolor
sollicitudin. Proin consequat magna quis lacus malesuada, sed tincidunt mauris
consequat. Integer posuere erat a ante venenatis dapibus posuere velit aliquet.
Nunc auctor, arcu vitae dignissim tincidunt, purus lectus suscipit metus, vel
vulputate nisl neque vitae arcu.

Pellentesque habitant morbi tristique senectus et netus et malesuada fames ac
turpis egestas. Quisque volutpat condimentum velit. Class aptent taciti sociosqu
ad litora torquent per conubia nostra, per inceptos himenaeos. In consectetur,
elit quis fermentum viverra, neque arcu pellentesque lectus, vitae tincidunt
libero justo ac ipsum. Aenean commodo ligula eget dolor. Aenean massa. Cum
sociis natoque penatibus et magnis dis parturient montes, nascetur ridiculus mus.
Donec quam felis, ultricies nec, pellentesque eu, pretium quis, sem.

\section{Template features}

\begin{itemize}
\item Changed font from default computer modern to CMBright.
\item Increased line spacing in the title to prevent letters from overlapping.
\item The siunitx package is included; see examples at: \url{https://github.com/dramco-edu/LaTex/wiki/LaTex-Tips-and-Tricks#using-si-units-numbers-angles-etc}
\item Include units alongside equations using the following command:
\begin{lstlisting}
    \begin{equation}
     E = m c^2
     \tagaddtext{[\si{\joule}]}
    \end{equation}
    \end{lstlisting},
which outputs:
\begin{equation}\label{eq:emc2}
    E = m c^2%
    \tagaddtext{[\si{\joule}]}
\end{equation}
\item Abbreviations:
\begin{itemize}
    \item  \gls{iot}      % outputs (if first occurance): Internet-of-Things (IoT)
          % outputs (else): IoT
    \item \glspl{lpwan}  % outputs (if first occurance): Low-Power Wide-Area Networks (LPWANs)
          % outputs (else): LPWANs
    \item \acrlong{iot}  % outputs: Internet-of-Things
    \item \acrshort{iot} % outputs: IoT
    \item \acrfull{iot}  % outputs: Internet-of-Things (IoT)
\end{itemize}
\item \verb!\ce{CO2}! \ce{CO2}
\item Refs:
\begin{itemize}

    \item \cref{fig:example}
    \item \cref{tab:example}
    \item \cref{ch:introduction}
    \item \cref{eq:emc2}


    \item \verb!\cref{fig:example}! becomes \cref{fig:example}
    \item \verb!\cref{tab:example}! becomes \cref{tab:example}
    \item \verb!\cref{ch:introduction}! becomes \cref{ch:introduction}
    \item \cref{eq:emc2}


\end{itemize}
\item numerical citation~\cite{ottoylorawan}
\begin{itemize}
    \item \verb!\cite{callebaut10lorawan}! becomes \cite{callebaut10lorawan}

    \item \verb!\citeauthor{callebaut10lorawan}! becomes \citeauthor{callebaut10lorawan}
    \item \verb!\citeauthorref{callebaut10lorawan}! becomes \citeauthorref{callebaut10lorawan}
    \item \verb!\autocite[chap.~2]{callebaut10lorawan}! becomes \autocite[chap.~2]{callebaut10lorawan}
    \item \verb!\footcite{callebaut10lorawan}! becomes \footcite{callebaut10lorawan}
    \item \verb!\fullcite{callebaut10lorawan}! becomes \fullcite{callebaut10lorawan}
\end{itemize}

    \item A $\SI{45}{\degree}$ angle or a \ang{45}.

    \item It is \SI{17}{\degreeCelsius} outside.

    \item  \num{10000}
    \item \num{3.45d-4}

    \item \si{\kilo\gram\meter\per\square\second} 
    \item \si{kg.m/s^2} %unit only

    \item \SI{10}{\percent}
    \item \SI{68}{kg}
\end{itemize}


\begin{table*}[t]
    \centering
    \begin{tabular}{ l  l  l  p{5cm}}
        \hline
        Day       & Min Temp & Max Temp & Summary                                                      \\ \hline
        Monday    & 11C      & 22C      & A clear day with lots of sunshine.
        However, the strong breeze will bring down the temperatures.                                   \\ \hline
        Tuesday   & 9C       & 19C      & Cloudy with rain, across many northern regions. Clear spells
        across most of Scotland and Northern Ireland,
        but rain reaching the far northwest.                                                           \\ \hline
        Wednesday & 10C      & 21C      & Rain will still linger for the morning.
        Conditions will improve by early afternoon and continue
        throughout the evening.                                                                        \\
        \hline
    \end{tabular}
\end{table*}







\newcolumntype{R}{>{\raggedleft\arraybackslash}X}%


\begin{table*}[t]
    \centering
    \begin{tabularx}{\textwidth}{ cccX }
        \hline
        Day       & Min Temp & Max Temp & Summary                                                      \\ \hline
        Monday    & 11C      & 22C      & A clear day with lots of sunshine.
        However, the strong breeze will bring down the temperatures.                                   \\ \hline
        Tuesday   & 9C       & 19C      & Cloudy with rain, across many northern regions. Clear spells
        across most of Scotland and Northern Ireland,
        but rain reaching the far northwest.                                                           \\ \hline
        Wednesday & 10C      & 21C      & Rain will still linger for the morning.
        Conditions will improve by early afternoon and continue
        throughout the evening.                                                                        \\
        \hline
    \end{tabularx}
\end{table*}



\begin{figure}
    \centering
    \includegraphics[width=0.95\linewidth]{example.jpg}
    \caption{Example JPG}%
    \label{fig:example}
\end{figure}

\begin{table*}[t]
    \centering
    \caption{Fixed-width columns.}%
    \label{tab:example}
    \begin{tabular}[t]{l>{\raggedright}p{0.3\linewidth}>{\raggedright\arraybackslash}p{0.3\linewidth}}
        \toprule
                     & Treatment A                    & Treatment B                     \\
        \midrule
        John Smith   & Good response, no side-effects & No response                     \\
        Jane Doe     & --                             & Good response, no side-effects  \\
        Mary Johnson & No response                    & Good response with side-effects \\
        \bottomrule
    \end{tabular}
\end{table*}%



\begin{listing*}[t]
    \inputminted[breaklines]{python}{code/example.py}
    \caption{Minimal working example}
    \label{listing:1}
\end{listing*}

\begin{figure}
    \centering
    % This file was created by tikzplotlib v0.9.8.
\begin{tikzpicture}



\begin{axis}[
width=\linewidth,
height=0.75\linewidth,
axis line style={white},
tick align=outside,
tick pos=left,
title={Simple plot \(\dfrac{\alpha}{2}\)},
xlabel={Time (\si{\second})},
xmajorgrids,
xmin=-0.095, xmax=1.995,
ylabel={Voltage (\si{\milli\volt})},
ymajorgrids,
ymin=-1.1, ymax=1.1,
]
\addplot [color0, mark=*]
table {%
0 0
0.1 0.587785252292473
0.2 0.951056516295154
0.3 0.951056516295154
0.4 0.587785252292473
0.5 1.22464679914735e-16
0.6 -0.587785252292473
0.7 -0.951056516295154
0.8 -0.951056516295154
0.9 -0.587785252292473
1 -2.44929359829471e-16
1.1 0.587785252292474
1.2 0.951056516295154
1.3 0.951056516295154
1.4 0.587785252292473
1.5 3.67394039744206e-16
1.6 -0.587785252292473
1.7 -0.951056516295154
1.8 -0.951056516295154
1.9 -0.587785252292473
};
\addlegendentry{sine}
\addplot [color1, mark=*]
table {%
0 1
0.1 0.809016994374947
0.2 0.309016994374947
0.3 -0.309016994374948
0.4 -0.809016994374947
0.5 -1
0.6 -0.809016994374947
0.7 -0.309016994374948
0.8 0.309016994374947
0.9 0.809016994374947
1 1
1.1 0.809016994374947
1.2 0.309016994374947
1.3 -0.309016994374947
1.4 -0.809016994374947
1.5 -1
1.6 -0.809016994374948
1.7 -0.309016994374946
1.8 0.309016994374947
1.9 0.809016994374947
};
\addlegendentry{cosine}
\end{axis}

\end{tikzpicture}

    \caption{Caption}
    \label{fig:my_label}
\end{figure}
